\section{Generalities}\label{sec:generalities}

\begin{itemize}
\item We write \(\R^{d + 1}_+ \eqd \R^d \times (0, \infty)\).
\item We write \(\langle x \rangle \eqd \sqrt{1 + | x |^2}\).
\item All \(L^p (\R^{d + 1}_+)\)-norms below are taken with respect to \(\dd x \frac{\dd t}{t}\).
\item We write \(\{\alpha\} \eqd \alpha - \lfloor \alpha \rfloor \in [0, 1)\).
\item A set \(\mathcal{B} \subset \Sch (\R^d)\) is bounded in \(\Sch (\R^d)\) if for every nonnegative integer \(M\),
\[
\sup_{\varphi \in \mathcal{B}} \sup_{z \in \R^d}
\sup_{| \beta | \leq M}
\langle z \rangle^M | \partial^{\beta} \varphi (z) | < \infty .
\]
\item The Fourier transform operator \(\Fourier\) is normalized by
\[
\FT{\varphi} (\xi) = (\Fourier \varphi) (\xi)
\eqd \int_{\R^d} \varphi (x) e^{- 2 \pi i x \cdot \xi} \dd x .
\]
\end{itemize}

\begin{definition}[Cancellation classes] \label{def:psi-classes} The set \(\Psi (\R^d) \subset \Sch (\R^d)\) is the set of Schwartz functions whose Fourier transform is supported away from the origin:
\begin{equation} \label{eq:psi-fourier-support}
\Psi (\R^d) \eqd \left\{ f \in \Sch (\R^d) \st 0 \not\in \spt \FT{f} \right\}.
\end{equation}
For a nonnegative integer \(m\), the set \(\Psi^m (\R^d) \subset \Sch (\R^d)\) is the set of Schwartz functions with vanishing moments up to order \(m\):
\begin{equation} \label{eq:psi-vanishing-moments}
\Psi^m (\R^d) \eqd \left\{ f \in \Sch (\R^d) \st
\int_{\R^d} f(x) x^{\beta}\dd x = 0
\quad \text{for every multi-index } \beta \text{ with } |\beta|\leq m\right\}.
\end{equation}
In particular, \(\Psi(\R^d) \subset \Psi^m(\R^d)\) for every nonnegative integer \(m\).
\end{definition}

\begin{definition}[H{\"o}lder spaces] \label{def:holder-spaces} Let \(\alpha > 0\) be noninteger and set \(n\eqd \lfloor \alpha \rfloor\).

For \(f \in C^n(\R^d)\), define the homogeneous H{\"o}lder seminorm by
\begin{equation}
\label{eq:holder-homogeneous-seminorm}
\|f\|_{\dot{\mathcal{C}}^{\alpha}(\R^d)} \eqd
\max_{|\gamma|=n} \sup_{x \neq x'}
\frac{| \partial^\gamma f(x') - \partial^\gamma f(x) |} {|x'-x|^{\{\alpha\}}}.
\end{equation}
This is a seminorm; it does not see polynomials of degree at most \(n\).

For \(t_0>0\), define the nonhomogeneous H{\"o}lder norm at top scale \(t_0\) by
\begin{equation} \label{eq:holder-nonhomogeneous-top-scale-norm}
\|f\|_{\mathcal{C}^{\alpha}_{t_0}(\R^d)} \eqd
\sum_{|\gamma|\le n} t_0^{|\gamma|} \|\partial^\gamma f\|_{L^\infty}
+
t_0^\alpha
\max_{|\gamma|=n}
\sup_{\substack{x\neq x' \\ |x'-x|<t_0 }}
\frac{| \partial^\gamma f(x')-\partial^\gamma f(x)|} {|x'-x|^{\{\alpha\}}}.
\end{equation}
\end{definition}

\begin{remark}[Homogeneity of the nonhomogeneous class] Although the nonhomogeneous H{\"o}lder norm is not scale-invariant for a fixed \(t_0\), the spaces are homogeneous as a class.
Specifically, spatial scaling simply shifts the top scale parameter.

More explicitly, setting \(f_\lambda(x) = f(x/\lambda)\) for \(\lambda > 0\), we obtain
\[
\|f_\lambda\|_{\mathcal{C}^{\alpha}_{\lambda t_0}(\R^d)}
=
\|f\|_{\mathcal{C}^{\alpha}_{t_0}(\R^d)}.
\]
\end{remark}

\begin{definition}[Wavelet operators] \label{def:wavelet-operators} For \(x \in \R^d\) and \(t > 0\), let the translation and dilation operators acting on a fixed test function \(\varphi \in \Sch(\R^d)\) be defined respectively by
\[
\Tr_x\varphi(y) \eqd \varphi(y-x)
\quad \text{and} \quad
\Dil_t\varphi(y) \eqd t^{-d}\varphi(y/t).
\]
\begin{enumerate}[(i)]
\item \textbf{Wavelet coefficient operator:} For a function \(f\) of at most polynomial growth, define
\begin{equation} \label{eq:wavelet-coefficient-fixed-test}
\Em_{\varphi}[f](x,t) \eqd
\int_{\R^d} f(y)\Tr_x\Dil_t\varphi(y)\dd y.
\end{equation}

\item \textbf{Continuous synthesis operator:} Let \(F:\R^{d+1}_+\to\C\).
Whenever the integral is absolutely convergent, define
\begin{equation} \label{eq:synthesis-operator}
(\dot{\mc{R}}_{\varphi}[F])(x) \eqd
\int_{\R^{d+1}_{+}} F(y,t)\Tr_y\Dil_t\varphi(x) \dd y\frac{\dd t}{t}.
\end{equation}

\item \textbf{Top-scale synthesis operator:} For \(t_0>0\) and a function \(a:\R^d\to\C\), define
\begin{equation} \label{eq:top-scale-synthesis-operator}
(\mc{R}^{\mrm{top}}_{\varphi}[a])(x,t_0) \eqd
\int_{\R^d} a(y)\Tr_y\Dil_{t_0}\varphi(x)\dd y.
\end{equation}
\end{enumerate}
\end{definition}

\begin{remark} \label{rem:variable-test-functions} One can more generally consider these operators when the test functions are allowed to depend on the space and scale parameters, provided they remain uniformly bounded in the Schwartz topology.
\end{remark}

\subsection{Wavelet coefficient characterization of H{\"o}lder spaces}

\begin{definition}[Calder\'on pairs]\label{defn:Calderon-pair} A homogeneous Calder\'on pair is a pair of functions \((\psi,\tilde{\psi})\in\Psi^{1}(\R^{d})\times \Psi^{1}(\R^{d})\) such that
\[
\int_{0}^{\infty}
\FT{\psi}(-t\xi)\FT{\tilde{\psi}} (t\xi)
\frac{\dd t}{t}
=
1
\qquad \text{for all } \xi\neq 0.
\]
The function \(\tilde{\psi}\) is called a Calder\'on conjugate for \(\psi\).

A pair
\[
((\psi,\varphi),(\tilde{\psi},\tilde{\varphi}))
\in
(\Psi^{1}(\R^{d})\times\Sch(\R^{d}))^{2}
\]
is a nonhomogeneous Calder\'on pair if
\[
\int_{0}^{1}
\FT{\psi}(-t\xi)\FT{\tilde{\psi}} (t\xi)
\frac{\dd t}{t}
+
\FT{\varphi}(-\xi)\FT{\tilde{\varphi}}(\xi)
=
1
\qquad \text{for all } \xi\in \R^d.
\]
The pair \((\tilde{\psi},\tilde{\varphi})\) is called a Calder\'on conjugate for \((\psi,\varphi)\).
\end{definition}

In \Cref{sec:Calderon-formula} we will see that, under mild assumptions, a wavelet \(\psi\in\Psi(\R^{d})\) has a conjugate \(\tilde{\psi}\) making \((\psi,\tilde{\psi})\in\Psi^{1}(\R^{d})\times \Psi^{1}(\R^{d})\) into a homogeneous Calder\'on pair.
Analogously, a pair \((\psi,\varphi)\in\Psi(\R^{d})\times\Sch(\R^{d})\), under mild assumptions, admits a conjugate \((\tilde{\psi},\tilde{\varphi})\in\Psi(\R^{d})\times\Sch(\R^{d})\), making \(((\psi,\varphi),(\tilde{\psi},\tilde{\varphi}))\in (\Psi^{1}(\R^{d})\times\Sch(\R^{d}))^{2}\) into a nonhomogeneous Calder\'on pair.
In particular, every open neighborhood of \(0\) in \(\Psi(\R^d)\), and every open neighborhood of \((0,0)\) in \(\Psi(\R^d)\times\Sch(\R^d)\), contains such wavelets.

\begin{theorem}[Wavelet coefficient characterization of Hölder spaces] \label{thm:holder-analysis} Let \(\alpha>0\) be noninteger.
For every \(\psi\in\Psi^{\lfloor\alpha\rfloor}(\R^d)\) and every \(f\in C^{\lfloor\alpha\rfloor}(\R^d)\) with \(\|f\|_{\dot{\mathcal{C}}^\alpha(\R^d)}<\infty\),
\begin{equation}
\label{eq:holder-analysis-homogeneous-leq}
\left\| t^{-\alpha}\Em_\psi[f]\right\|_{L^\infty(\R^{d+1}_+)}
\lesssim
\|f\|_{\dot{\mathcal{C}}^\alpha(\R^d)},
\end{equation}
with an implicit constant uniformly bounded when \(\psi\in\Psi^{\lfloor\alpha\rfloor}(\R^{d})\) varies in a bounded set.

Furthermore, if \(\psi\) admits a Calder\'on conjugate, then
\begin{equation}
\label{eq:holder-analysis-homogeneous-geq}
\left\| t^{-\alpha}\Em_\psi[f]\right\|_{L^\infty(\R^{d+1}_+)}
\gtrsim
\|f\|_{\dot{\mathcal{C}}^\alpha(\R^d)}.
\end{equation}
In particular, for any open neighborhood \(\mf{U}\) of \(0\) in \(\Psi(\R^{d})\), there exists \(\psi\in\mf{U}\) such that the bound \eqref{eq:holder-analysis-homogeneous-geq} holds.

For every \((\psi,\varphi)\in\Psi^n(\R^d)\times\Sch(\R^d)\), every \(t_0>0\), and every \(f\in\mathcal{C}^{\alpha}_{t_0}(\R^d)\),
\begin{equation}
\label{eq:holder-analysis-nonhomogeneous-top-scale-leq}
\begin{aligned}[t]
   & \left\|\Em_\varphi[f](\cdot,t_0)\right\|_{L^\infty(\R^d)}
  + t_0^\alpha
  \left\| \1_{0<t<t_0}t^{-\alpha}\Em_\psi[f] \right\|_{L^\infty(\R^{d+1}_+)}
  \\
   & \qquad\lesssim \|f\|_{\mathcal{C}^{\alpha}_{t_0}(\R^d)},
\end{aligned}
\end{equation}
with an implicit constant uniformly bounded when \((\psi,\varphi)\in\Psi^{\lfloor\alpha\rfloor}(\R^{d})\times\Sch(\R^{d})\) varies in a bounded set.

Furthermore, if \((\psi,\varphi)\) admits a Calder\'on conjugate, then
\begin{equation}
\label{eq:holder-analysis-nonhomogeneous-top-scale-geq}
\begin{aligned}[t]
   &
  \left\|\Em_\varphi[f](\cdot,t_0)\right\|_{L^\infty(\R^d)}
  +
  t_0^\alpha
  \left\|
  \1_{0<t<t_0}t^{-\alpha}\Em_\psi[f]
  \right\|_{L^\infty(\R^{d+1}_+)}
  \gtrsim
  \|f\|_{\mathcal{C}^{\alpha}_{t_0}(\R^d)}.
\end{aligned}
\end{equation}
In particular, for any open neighborhood \(\mf{U}\) of \((0,0)\) in \(\Psi(\R^{d})\times\Sch(\R^{d})\), there exists \((\psi,\varphi)\in\mf{U}\) such that the bound \eqref{eq:holder-analysis-nonhomogeneous-top-scale-geq} holds.
\end{theorem}


The proof of upper bounds \eqref{eq:holder-analysis-homogeneous-leq} and \eqref{eq:holder-analysis-nonhomogeneous-top-scale-leq} follows by direct computation as we show below. The lower bounds, rely on being able to reconstruct the orginal function \(f\) from its coefficiens \(\mc{E}[f]\), as stated in the following theorem. 

\begin{theorem}[Calder{\'o}n reproducing formulas] \label{thm:calderon-reproducing} Let \((\psi,\tilde{\psi})\in\Psi^{1}(\R^d)\times\Psi^{1}(\R^d)\) be a homogeneous Calder\'on pair.
Then, for every \(f\in\Sch(\R^d)\),
\begin{equation}
\label{eq:homogeneous-calderon-formula}
\lim_{\varepsilon\to0^+}
\dot{\mc{R}}_{\tilde{\psi}}
\left[\1_{\varepsilon<t<\varepsilon^{-1}}\Em_\psi[f]\right]
=
f
\qquad \text{in } \Sch'(\R^d).
\end{equation}
The same identity holds for every \(f\in\Sch'(\R^d)\) modulo polynomials, with coefficients interpreted by distributional pairing:
\[
\Em_\psi[f](x,t)\eqd
\left\langle f,\Tr_x\Dil_t\psi\right\rangle .
\]

Let \(((\psi,\varphi),(\tilde{\psi},\tilde{\varphi})) \in (\Psi^{1}(\R^d)\times\Sch(\R^d))^{2}\) be a nonhomogeneous Calder\'on pair.
Then, for every \(t_0>0\) and every \(f\in\Sch(\R^d)\),
\begin{equation}
\label{eq:nonhomogeneous-calderon-formula-top-scale}
\lim_{\varepsilon\to0^+}
\dot{\mc{R}}_{\tilde{\psi}}
\left[\1_{\varepsilon<t<t_0}\Em_\psi[f]\right]
+
\left(
\mc{R}^{\mrm{top}}_{\tilde{\varphi}}
\left[\Em_\varphi[f](\cdot,t_0)\right]
\right)(\cdot,t_0)
=
f
\qquad \text{in } \Sch'(\R^d).
\end{equation}
The same identity holds for every \(f\in\Sch'(\R^d)\), with coefficients interpreted by distributional pairing.
\end{theorem}
We proved the proof of this theorem, as well as a method for constructing Calderón conjugates in \Cref{sec:Calderon-formula}, below. 



\begin{proof}[Proof of upper bounds \eqref{eq:holder-analysis-homogeneous-leq} and \eqref{eq:holder-analysis-nonhomogeneous-top-scale-leq}] The homogeneous estimate is the usual Taylor estimate with cancellation.
Indeed, after \(y=x+tz\),
\[
\Em_\psi[f](x,t)=\int_{\R^d} f(x+tz)\psi(z)\dd z.
\]
Let \(P_{x,n}\) denote the Taylor polynomial of \(f\) at \(x\) of degree \(n\).
Since the moments of \(\psi\) up to order \(n\) vanish,
\[
\int_{\R^d} P_{x,n}(x+tz)\psi(z)\dd z=0.
\]
Therefore,
\[
|\Em_\psi[f](x,t)|
\le
\int_{\R^d}|f(x+tz)-P_{x,n}(x+tz)|\,|\psi(z)|\dd z
\lesssim_{d,\alpha}
t^\alpha\|f\|_{\dot{\mathcal{C}}^\alpha}
\int_{\R^d}|z|^\alpha|\psi(z)|\dd z.
\]
The last integral is uniformly bounded when \(\psi\) varies in a bounded subset of \(\Sch(\R^d)\), which gives \eqref{eq:holder-analysis-homogeneous-leq}.

Let \(A\eqd \|f\|_{\mathcal{C}^{\alpha}_{t_0}}\).
If \(0<t<t_0\), the same Taylor argument gives
\[
|f(x+tz)-P_{x,n}(x+tz)|
\lesssim_{d,\alpha}
A \left(\frac{t}{t_0}\right)^\alpha
\langle z\rangle^\alpha .
\]
For \(|tz|<t_0\), this follows from the top-scale H{\"o}lder seminorm.
For \(|tz|\ge t_0\), it follows from the scaled bounds \(\|\partial^\gamma f\|_\infty\le A t_0^{-|\gamma|}\), since
\[
\left(\frac{t|z|}{t_0}\right)^j
\le
\left(\frac{t|z|}{t_0}\right)^\alpha
\qquad
\text{for } 0\le j\le n
\]
whenever \(t|z|/t_0\ge1\).
Hence
\[
|\Em_\psi[f](x,t)|
\lesssim_{d,\alpha,\psi}
A t_0^{-\alpha} t^\alpha ,
\qquad 0<t<t_0.
\]
Uniformly for \(\psi\) in a bounded subset of \(\Sch(\R^d)\), this gives the high-frequency term in \eqref{eq:holder-analysis-nonhomogeneous-top-scale-leq}.

For the top-scale term, no cancellation is required:
\[
|\Em_\varphi[f](x,t_0)|
\le
\|f\|_{L^\infty}\int_{\R^d}|\varphi(z)|\dd z
\lesssim_{\varphi} A.
\]
Uniformly for \(\varphi\) in a bounded subset of \(\Sch(\R^d)\), this gives the top-scale term in \eqref{eq:holder-analysis-nonhomogeneous-top-scale-leq}.
\end{proof}

\subsection{Calder{\'o}n formulas}\label{sec:Calderon-formula}
In this section we prove \Cref{thm:calderon-reproducing} and then describe the construction of Calderón conjugates, defined in \Cref{defn:Calderon-pair}.

\begin{proof} For \(f\in\Sch(\R^d)\), the definitions of \(\Em_\psi\) and \(\dot{\mc{R}}_{\tilde{\psi}}\) give
\[
\FT{\left(\Em_\psi[f](\cdot,t)\right)}(\xi)
=
\Fourier{f}(\xi)\FT{\psi}(-t\xi),
\]
and
\[
\Fourier{\left( \dot{\mc{R}}_{\tilde{\psi}} \left[\1_{\varepsilon<t<R}\Em_\psi[f]\right] \right)}(\xi)
=
\FT{f}(\xi) \int_{\varepsilon}^{R} \FT{\psi}(-t\xi)\FT{\tilde{\psi}}(t\xi) \frac{\dd t}{t}.
\]
Taking \(R=\varepsilon^{-1}\) and using the homogeneous Calder\'on identity, the multiplier converges pointwise to \(1\) for \(\xi\neq0\).
The convergence against Schwartz test functions follows from dominated convergence, using the first-order vanishing at the origin of the Fourier transforms of functions in \(\Psi^1(\R^d)\) and rapid decay at infinity.
This proves \eqref{eq:homogeneous-calderon-formula}.

For the nonhomogeneous formula, the top-scale term has Fourier transform
\[
\Fourier{\left(\left(\mc{R}^{\mrm{top}}_{\tilde{\varphi}} \left[\Em_\varphi[f](\cdot,t_0)\right] \right)(\cdot,t_0)\right)}(\xi)
=
\FT{f}(\xi) \FT{\varphi}(-t_0\xi)\FT{\tilde{\varphi}}(t_0\xi).
\]
Thus the multiplier associated with the left side of \eqref{eq:nonhomogeneous-calderon-formula-top-scale} is
\[
m_{\varepsilon,t_0}(\xi)
\eqd
\int_{\varepsilon}^{t_0} \FT{\psi}(-t\xi)\FT{\tilde{\psi}}(t\xi) \frac{\dd t}{t} + \FT{\varphi}(-t_0\xi)\FT{\tilde{\varphi}}(t_0\xi).
\]
By applying the nonhomogeneous Calder\'on identity at \(t_0\xi\), and changing variables \(t=s t_0\), we obtain
\[
\int_{0}^{t_0} \FT{\psi}(-t\xi)\FT{\tilde{\psi}}(t\xi) \frac{\dd t}{t} + \FT{\varphi}(-t_0\xi)\FT{\tilde{\varphi}}(t_0\xi) = 1.
\]
Hence \(m_{\varepsilon,t_0}(\xi)\to1\) for every \(\xi\in\R^d\), and again dominated convergence gives \eqref{eq:nonhomogeneous-calderon-formula-top-scale} in \(\Sch'(\R^d)\).

The extension to tempered distributions follows by transposition.
In the homogeneous case, transposition gives convergence on the quotient of \(\Sch'(\R^d)\) by the distributions whose Fourier transforms are supported at the origin, equivalently modulo polynomials.
\end{proof}

\begin{lemma}[Construction of Calder{\'o}n conjugates] \label{lem:calderon-conjugate-construction} Let \(\psi\in\Psi(\R^d)\) satisfy the ray non-degeneracy condition
\[
\text{for every } \xi\neq0,\quad
\text{there exists } s>0 \text{ such that } \FT{\psi}(s\xi)\neq0.
\]
Then there exists \(\tilde{\psi}\in\Psi(\R^d)\), with \(\FT{\tilde{\psi}}\in C^\infty_c(\R^d\setminus\{0\})\), such that
\[
\int_0^\infty \FT{\psi}(-t\xi)\FT{\tilde{\psi}}(t\xi)\frac{\dd t}{t} = 1, \qquad \xi\neq0.
\]
In particular, \((\psi,\tilde{\psi})\) is a homogeneous Calder\'on pair.

If, in addition, \(\varphi\in\Sch(\R^d)\) satisfies
\[
\FT{\varphi}(\xi)\neq0 \qquad \text{for every } \xi\in\R^d,
\]
then there exists \(\tilde{\varphi}\in\Sch(\R^d)\), with \(\FT{\tilde{\varphi}}\in C^\infty_c(\R^d)\), such that \(((\psi,\varphi),(\tilde{\psi},\tilde{\varphi}))
\) is a nonhomogeneous Calder\'on pair.

Consequently, every open neighborhood of \(0\) in \(\Psi(\R^d)\) contains a wavelet admitting a homogeneous Calder\'on conjugate, and every open neighborhood of \((0,0)\) in \(\Psi(\R^d)\times\Sch(\R^d)\) contains a pair admitting a nonhomogeneous Calder\'on conjugate.
\end{lemma}

\begin{proof}[Proof of \Cref{lem:calderon-conjugate-construction}] Set
\[
a(\eta)\eqd \FT{\psi}(-\eta).
\]
By applying the ray non-degeneracy assumption to \(-\omega\), for every \(\omega\in\mathbb{S}^{d-1}\) there exists \(r_\omega>0\) such that
\[
a(r_\omega\omega)\neq0.
\]
By compactness of \(\mathbb{S}^{d-1}\), choose finitely many open sets \(U_j\subset\mathbb{S}^{d-1}\), radii localized by nonnegative functions \(\rho_j\in C^\infty_c((0,\infty))\), and a smooth partition of unity \(\{\theta_j\}\) on \(\mathbb{S}^{d-1}\) subordinate to the \(U_j\), such that
\[
\int_0^\infty \rho_j(r)|a(r\omega)|^2\frac{\dd r}{r}>0
\qquad
\text{for every } \omega\in\spt \theta_j.
\]
For \(\eta=r\omega\), \(r>0\), define
\[
h(r\omega) \eqd \sum_j \theta_j(\omega)\rho_j(r)\overline{a(r\omega)}.
\]
Then \(h\in C^\infty_c(\R^d\setminus\{0\})\).
Define
\[
A(\omega) \eqd \int_0^\infty a(r\omega)h(r\omega)\frac{\dd r}{r}.
\]
By construction,
\[
A(\omega) = \sum_j \theta_j(\omega) \int_0^\infty \rho_j(r)|a(r\omega)|^2\frac{\dd r}{r} > 0 \qquad \text{for every } \omega\in\mathbb{S}^{d-1}.
\]
Now set
\[
\FT{\tilde{\psi}}(\eta) \eqd h(\eta)A(\eta/|\eta|)^{-1}, \qquad \eta\neq0,
\]
and extend by \(0\) at the origin.
Since \(h\) is compactly supported away from the origin, \(\FT{\tilde{\psi}}\in C^\infty_c(\R^d\setminus\{0\})\), hence \(\tilde{\psi}\in\Psi(\R^d)\).
For \(\xi=|\xi|\omega\neq0\),
\[
\begin{aligned}[t]
\int_0^\infty \FT{\psi}(-t\xi)\FT{\tilde{\psi}}(t\xi)\frac{\dd t}{t} & = A(\omega)^{-1} \int_0^\infty a(t|\xi|\omega)h(t|\xi|\omega)\frac{\dd t}{t}
\\
& = A(\omega)^{-1} \int_0^\infty a(r\omega)h(r\omega)\frac{\dd r}{r} = 1.
\end{aligned}
\]

For the nonhomogeneous construction, define
\[
g(\eta)\eqd \FT{\psi}(-\eta)\FT{\tilde{\psi}}(\eta).
\]
Since \(\FT{\tilde{\psi}}\in C^\infty_c(\R^d\setminus\{0\})\), we have \(g\in C^\infty_c(\R^d\setminus\{0\})\), and the homogeneous identity gives
\[
\int_0^\infty g(t\xi)\frac{\dd t}{t}=1 \qquad \text{for } \xi\neq0.
\]
Set
\[
b(\eta) \eqd 1-\int_0^1 g(t\eta)\frac{\dd t}{t}.
\]
Then \(b\in C^\infty_c(\R^d)\), \(b=1\) in a neighborhood of the origin, and
\[
b(\eta) = 1-\int_0^1 \FT{\psi}(-t\eta)\FT{\tilde{\psi}}(t\eta)\frac{\dd t}{t}.
\]
Define
\[
\FT{\tilde{\varphi}}(\eta) \eqd \frac{b(\eta)}{\FT{\varphi}(-\eta)}.
\]
The denominator is smooth and nonvanishing, while \(b\) is compactly supported and smooth; hence \(\FT{\tilde{\varphi}}\in C^\infty_c(\R^d)\), so \(\tilde{\varphi}\in\Sch(\R^d)\).
Moreover,
\[
\int_0^1 \FT{\psi}(-t\xi)\FT{\tilde{\psi}}(t\xi) \frac{\dd t}{t} + \FT{\varphi}(-\xi)\FT{\tilde{\varphi}}(\xi) = 1 \qquad \text{for every } \xi\in\R^d.
\]
Thus
\[
((\psi,\varphi),(\tilde{\psi},\tilde{\varphi}))
\]
is a nonhomogeneous Calder\'on pair.

Finally, choose a fixed \(\psi_0\in\Psi(\R^d)\) with \(\FT{\psi_0}\in C^\infty_c(\R^d\setminus\{0\})\) nonzero on an annulus meeting every ray.
Then \(\lambda\psi_0\) satisfies the ray non-degeneracy condition for every \(\lambda\neq0\), and \(\lambda\psi_0\to0\) in \(\Sch(\R^d)\) as \(\lambda\to0\).
This proves the homogeneous neighborhood assertion.
For the nonhomogeneous assertion, take in addition a fixed Gaussian \(\varphi_0\in\Sch(\R^d)\) with \(\FT{\varphi_0}>0\) everywhere, and use \((\lambda\psi_0,\mu\varphi_0)\) with nonzero \(\lambda,\mu\) sufficiently small.
\end{proof}

\subsection{Proof of upper bounds and synthesis estimates}

\begin{proof}[Proof of lower bounds \eqref{eq:holder-analysis-homogeneous-geq} and \eqref{eq:holder-analysis-nonhomogeneous-top-scale-geq}] Set \(n\eqd\lfloor\alpha\rfloor\) and \(\beta\eqd\{\alpha\}=\alpha-n\).

First suppose that \(\psi\) admits a homogeneous Calder\'on conjugate \(\tilde{\psi}\), and set
\[
A\eqd
\left\|t^{-\alpha}\Em_\psi[f]\right\|_{L^\infty(\R^{d+1}_+)}.
\]
For \(0<\varepsilon<1\), define
\[
G_\varepsilon
\eqd
\dot{\mc{R}}_{\tilde{\psi}}
\left[\1_{\varepsilon<t<\varepsilon^{-1}}\Em_\psi[f]\right].
\]
Write \(\Delta_u g(x)\eqd g(x+u)-g(x)\).
For every multi-index \(\gamma\) with \(|\gamma|=n\), every \(u\neq0\), and every \(\varepsilon\),
\[
\begin{aligned}
  \|\Delta_u\partial^\gamma G_\varepsilon\|_{L^\infty}
   & \le
  A
  \int_0^\infty
  t^\alpha
  \|\Delta_u\partial^\gamma\Dil_t\tilde{\psi}\|_{L^1}
  \frac{\dd t}{t}
  \\
   & \lesssim_{\tilde{\psi}}
  A
  \left(
  \int_0^{|u|} t^{\alpha-n}\frac{\dd t}{t}
  +
  |u|\int_{|u|}^{\infty} t^{\alpha-n-1}\frac{\dd t}{t}
  \right)
  \\
   & \lesssim_{\tilde{\psi},\alpha}
  A |u|^{\alpha-n}.
\end{aligned}
\]
Indeed,
\[
\|\Delta_u\partial^\gamma\Dil_t\tilde{\psi}\|_{L^1}
\lesssim_{\tilde{\psi}}
\min\{t^{-n},\, |u|t^{-n-1}\}.
\]
Passing to the limit in the homogeneous Calder\'on formula shows that \(f\) differs, modulo a polynomial, from a distribution whose \(\dot{\mathcal{C}}^\alpha\)-seminorm is bounded by \(A\).
Since both \(f\) and this representative have finite \(\dot{\mathcal{C}}^\alpha\)-seminorm, the polynomial ambiguity has degree at most \(n\), and hence is invisible to \(\|\cdot\|_{\dot{\mathcal{C}}^\alpha}\).
Therefore
\[
\|f\|_{\dot{\mathcal{C}}^\alpha(\R^d)}
\lesssim_{\psi,\tilde{\psi},\alpha}
A,
\]
which is \eqref{eq:holder-analysis-homogeneous-geq}.

Now suppose that \((\psi,\varphi)\) admits a nonhomogeneous Calder\'on conjugate \((\tilde{\psi},\tilde{\varphi})\), and set
\[
B\eqd
\left\|\Em_\varphi[f](\cdot,t_0)\right\|_{L^\infty(\R^d)}
+
t_0^\alpha
\left\|
\1_{0<t<t_0}t^{-\alpha}\Em_\psi[f]
\right\|_{L^\infty(\R^{d+1}_+)}.
\]
By \eqref{eq:nonhomogeneous-calderon-formula-top-scale},
\[
f=G+H
\qquad \text{in } \Sch'(\R^d),
\]
where
\[
G
\eqd
\dot{\mc{R}}_{\tilde{\psi}}
\left[\1_{0<t<t_0}\Em_\psi[f]\right],
\]
and
\[
H(\cdot)
\eqd
\left(
\mc{R}^{\mrm{top}}_{\tilde{\varphi}}
\left[\Em_\varphi[f](\cdot,t_0)\right]
\right)(\cdot,t_0).
\]
For every \(|\gamma|\le n\),
\[
\begin{aligned}
  \|\partial^\gamma G\|_{L^\infty}
   & \le
  B t_0^{-\alpha}
  \int_0^{t_0}
  t^\alpha
  \|\partial^\gamma\Dil_t\tilde{\psi}\|_{L^1}
  \frac{\dd t}{t}
  \\
   & \lesssim_{\tilde{\psi},\alpha}
  B t_0^{-|\gamma|},
\end{aligned}
\]
and
\[
\|\partial^\gamma H\|_{L^\infty}
\le
B\|\partial^\gamma\Dil_{t_0}\tilde{\varphi}\|_{L^1}
\lesssim_{\tilde{\varphi}}
B t_0^{-|\gamma|}.
\]
Thus the \(L^\infty\)-part of \(\|f\|_{\mathcal{C}^\alpha_{t_0}}\) is bounded by \(B\).

It remains to estimate the top-scale H{\"o}lder seminorm.
Let \(|\gamma|=n\) and \(0<|u|<t_0\).
As above,
\[
\begin{aligned}
  \|\Delta_u\partial^\gamma G\|_{L^\infty}
   & \le
  B t_0^{-\alpha}
  \int_0^{t_0}
  t^\alpha
  \|\Delta_u\partial^\gamma\Dil_t\tilde{\psi}\|_{L^1}
  \frac{\dd t}{t}
  \\
   & \lesssim_{\tilde{\psi},\alpha}
  B t_0^{-\alpha}|u|^{\alpha-n}.
\end{aligned}
\]
For the top-scale piece,
\[
\|\Delta_u\partial^\gamma H\|_{L^\infty}
\le
B\|\Delta_u\partial^\gamma\Dil_{t_0}\tilde{\varphi}\|_{L^1}
\lesssim_{\tilde{\varphi}}
B |u| t_0^{-n-1}.
\]
Since \(0<|u|<t_0\) and \(0<\beta<1\),
\[
|u|t_0^{-n-1}
\le
t_0^{-\alpha}|u|^{\alpha-n}.
\]
Therefore
\[
\max_{|\gamma|=n}
\sup_{\substack{x\neq x'\\ |x'-x|<t_0}}
\frac{
|\partial^\gamma f(x')-\partial^\gamma f(x)|
}{|x'-x|^{\{\alpha\}}}
\lesssim_{\tilde{\psi},\tilde{\varphi},\alpha}
B t_0^{-\alpha}.
\]
Combining the derivative bounds and the H{\"o}lder bound gives
\[
\|f\|_{\mathcal{C}^{\alpha}_{t_0}(\R^d)}
\lesssim_{\psi,\varphi,\tilde{\psi},\tilde{\varphi},\alpha}
B,
\]
which is exactly \eqref{eq:holder-analysis-nonhomogeneous-top-scale-geq}.
\end{proof}


\begin{lemma}[Wave packet kernel bounds] \label{lem:wave-packet-kernel-bounds} Let \(φ∈\Sch(\R^{d})\). For every multi-index \(\gamma\),
\begin{equation}
\label{eq:WP-derivative-and-increment}
\begin{aligned}
   & \int_{\R^d}
  \left|
  t^{|\gamma|}
  \partial_y^\gamma\Tr_x\Dil_tφ(y)
  \right|\dd x
  \lesssim 1,    \\
   & \int_{\R^d}
  t^{|\gamma|}
  \left|
  \partial_y^\gamma\Tr_x\Dil_tφ(y')
  -
  \partial_y^\gamma\Tr_x\Dil_tφ(y)
  \right|\dd x
  \lesssim
  \min\left(1,\frac{|y'-y|}{t}\right).
\end{aligned}
\end{equation}
The implicit constants depend only on finitely many Schwartz seminorms of \(φ\).
\end{lemma}

\begin{proof} The first estimate follows from
\[
\partial_y^\gamma\Tr_x\Dil_tφ(y)
=
t^{-d-|\gamma|}
(\partial^\gammaφ)((y-x)/t)
\]
and the change of variables \(z=(y-x)/t\), using uniform Schwartz bounds.
For the increment estimate, when \(|y'-y|\le t\), apply the mean value theorem and the first estimate with \(|\gamma|+1\) in place of \(|\gamma|\).
When \(|y'-y|>t\), use the first estimate on each term separately.
\end{proof}

\begin{proposition}[Homogeneous and nonhomogeneous synthesis estimates] \label{prop:holder-synthesis} Let \(\alpha>0\) be noninteger, set
\[
n\eqd \lfloor\alpha\rfloor,
\qquad
\delta\eqd \alpha-n\in(0,1).
\]
We use the convention
\[
(\Tr_x\Dil_t\psi)(y)=t^{-d}\psi\!\left(\frac{y-x}{t}\right).
\]

\smallskip \noindent \emph{Homogeneous estimate.} Let \(F:\R^{d+1}_+\to\C\) be measurable and assume
\[
A_\infty\eqd
\left\|t^{-\alpha}F\right\|_{L^\infty(\R^{d+1}_+)}
<\infty .
\]
Let \(\Phi=(\varphi_{x,t})_{(x,t)\in\R^{d+1}_+}\) be bounded in \(\Sch(\R^d)\).
For each \((x,t)\in\R^{d+1}_+\), define
\[
\varphi^{\#}_{x,t}(z)
\eqd
\varphi_{x,t}(z)
-
\sum_{|\eta|\le n}
\frac{\left(z+\frac{x}{t}\right)^\eta}{\eta!}
\left(\partial^\eta\varphi_{x,t}\right)\!\left(-\frac{x}{t}\right).
\]
Equivalently,
\[
(\Tr_x\Dil_t\varphi^{\#}_{x,t})(y)
=
(\Tr_x\Dil_t\varphi_{x,t})(y)
-
\sum_{|\eta|\le n}
\frac{y^\eta}{\eta!}
\partial_y^\eta
\left(\Tr_x\Dil_t\varphi_{x,t}\right)(0).
\]
Set
\[
\left(\dot{\mc{R}}^0_{\Phi}[F]\right)(y)
\eqd
\int_{\R^{d+1}_+}
F(x,t)
(\Tr_x\Dil_t\varphi^{\#}_{x,t})(y)
\dd x\frac{\dd t}{t}.
\]
Then the integral converges absolutely for every \(y\), locally uniformly after applying \(\partial^\beta\) for every \(|\beta|\le n\), and
\[
\left\|\dot{\mc{R}}^0_{\Phi}[F]\right\|_{\dot{\mathcal{C}}^\alpha(\R^d)}
\lesssim
A_\infty .
\]

\smallskip \noindent \emph{Nonhomogeneous estimate at top scale.} Let \(t_0>0\).
Let
\[
\Phi=(\varphi_{x,t})_{(x,t)\in\R^d\times(0,t_0)}
\quad\text{and}\quad
\Phi^0=(\varphi^0_x)_{x\in\R^d}
\]
be bounded families in \(\Sch(\R^d)\).
Let \(F_0:\R^d\times(0,t_0)\to\C\) and \(F_1:\R^d\to\C\) be measurable, and assume
\[
A_{t_0}\eqd
t_0^\alpha
\left\|
\1_{0<t<t_0}t^{-\alpha}F_0
\right\|_{L^\infty(\R^{d+1}_+)}
+
\|F_1\|_{L^\infty(\R^d)}
<\infty .
\]
Define
\begin{equation}
\label{eq:nonhomogeneous-synthesis-definition-top-scale}
G(y)
\eqd
\int_0^{t_0}\int_{\R^d}
F_0(x,t)
(\Tr_x\Dil_t\varphi_{x,t})(y)
\dd x\frac{\dd t}{t}
+
\int_{\R^d}
F_1(x)
(\Tr_x\Dil_{t_0}\varphi^0_x)(y)
\dd x .
\end{equation}
Then \(G\in\mathcal{C}^{\alpha}_{t_0}(\R^d)\), the integrals in \eqref{eq:nonhomogeneous-synthesis-definition-top-scale} converge absolutely pointwise and locally uniformly after applying \(\partial^\beta\) for every \(|\beta|\le n\), and
\begin{equation}
\label{eq:holder-synthesis-nonhomogeneous-top-scale}
\|G\|_{\mathcal{C}^{\alpha}_{t_0}(\R^d)}
\lesssim
A_{t_0}.
\end{equation}

The implicit constants depend only on \(d,\alpha\), and finitely many Schwartz seminorms of the relevant families.
In the nonhomogeneous estimate, the implicit constant is independent of \(t_0\).
\end{proposition}

\begin{proof} We first record the common wave-packet bounds used in both estimates.
Let \(\Psi=(\psi_{x,t})\) be any bounded family in \(\Sch(\R^d)\).
Then, for every multiindex \(\theta\),
\begin{equation}
\label{eq:common-wave-packet-l1-bound}
\sup_{y\in\R^d}
\int_{\R^d}
\left|
\partial_y^\theta
(\Tr_x\Dil_t\psi_{x,t})(y)
\right|
\dd x
\lesssim
t^{-|\theta|}.
\end{equation}
Indeed, choosing \(N>d\), boundedness of \(\Psi\) in \(\Sch(\R^d)\) gives
\[
\left|
\partial_y^\theta
(\Tr_x\Dil_t\psi_{x,t})(y)
\right|
\lesssim
t^{-d-|\theta|}
\left(1+\frac{|y-x|}{t}\right)^{-N},
\]
and integration in \(x\) gives \eqref{eq:common-wave-packet-l1-bound}.
Consequently, for \(h=|y'-y|\),
\begin{equation}
\label{eq:common-wave-packet-difference-bound}
\int_{\R^d}
\left|
\partial_y^\theta
(\Tr_x\Dil_t\psi_{x,t})(y')
-
\partial_y^\theta
(\Tr_x\Dil_t\psi_{x,t})(y)
\right|
\dd x
\lesssim
t^{-|\theta|}
\min\left(1,\frac{h}{t}\right).
\end{equation}
The first bound in the minimum follows from \eqref{eq:common-wave-packet-l1-bound}, while the second follows from the mean value theorem and \eqref{eq:common-wave-packet-l1-bound} applied with \(|\theta|+1\) derivatives.
The same estimates apply to a bounded family \((\psi_x)_{x\in\R^d}\) at a fixed scale \(t\), by taking \(\psi_{x,t}=\psi_x\).

We prove the homogeneous estimate first.
Since
\[
|F(x,t)|\le A_\infty t^\alpha,
\]
fix a compact set \(E\subset\R^d\) and \(|\beta|\le n\).
Write
\[
g_{x,t}(y)\eqd(\Tr_x\Dil_t\varphi_{x,t})(y),
\qquad
g^\#_{x,t}(y)\eqd(\Tr_x\Dil_t\varphi^\#_{x,t})(y).
\]
By the defining identity for \(g^\#_{x,t}\),
\[
\partial_y^\beta g^\#_{x,t}(y)
=
\partial_y^\beta g_{x,t}(y)
-
\sum_{\substack{|\eta|\le n\\
\beta\le \eta}}
\frac{y^{\eta-\beta}}{(\eta-\beta)!}
\partial_y^\eta g_{x,t}(0).
\]
Using \eqref{eq:common-wave-packet-l1-bound}, uniformly for \(y\in E\),
\[
\int_{\R^d}
\left|
\partial_y^\beta g^\#_{x,t}(y)
\right|
\dd x
\lesssim_E
\sum_{|\eta|\le n} t^{-|\eta|}
\lesssim_E
t^{-n},
\qquad 0<t\le 1.
\]
Hence
\[
\sup_{y\in E}
\int_0^1\int_{\R^d}
|F(x,t)|
\left|
\partial_y^\beta g^\#_{x,t}(y)
\right|
\dd x\frac{\dd t}{t}
\lesssim_E
A_\infty
\int_0^1 t^{\alpha-n}\frac{\dd t}{t}
<\infty,
\]
because \(\alpha>n\).

For \(t\ge 1\), Taylor's theorem with integral remainder, applied to \(\partial_y^\beta g_{x,t}\) at \(0\), gives, uniformly for \(y\in E\),
\[
\left|
\partial_y^\beta g^\#_{x,t}(y)
\right|
\lesssim_E
\sum_{|\theta|=n+1}
\int_0^1
\left|
\partial_y^\theta g_{x,t}(sy)
\right|
\dd s .
\]
Therefore, by \eqref{eq:common-wave-packet-l1-bound},
\[
\sup_{y\in E}
\int_{\R^d}
\left|
\partial_y^\beta g^\#_{x,t}(y)
\right|
\dd x
\lesssim_E
t^{-n-1},
\qquad t\ge 1.
\]
Thus
\[
\sup_{y\in E}
\int_1^\infty\int_{\R^d}
|F(x,t)|
\left|
\partial_y^\beta g^\#_{x,t}(y)
\right|
\dd x\frac{\dd t}{t}
\lesssim_E
A_\infty
\int_1^\infty t^{\alpha-n-1}\frac{\dd t}{t}
<\infty,
\]
because \(\alpha<n+1\).
This proves the asserted absolute convergence and locally uniform convergence after applying \(\partial^\beta\) for every \(|\beta|\le n\).

It remains, in the homogeneous case, to prove the Hölder estimate.
Let \(|\gamma|=n\), \(y,y'\in\R^d\), and \(h=|y'-y|>0\).
Since the correction term in \(g^\#_{x,t}\) is a polynomial in \(y\) of degree at most \(n\), its \(\gamma\)-derivative is constant.
Hence
\[
\partial_y^\gamma g^\#_{x,t}(y')
-
\partial_y^\gamma g^\#_{x,t}(y)
=
\partial_y^\gamma g_{x,t}(y')
-
\partial_y^\gamma g_{x,t}(y).
\]
Using \eqref{eq:common-wave-packet-difference-bound},
\[
\begin{aligned}
   &
  \left|
  \partial^\gamma\dot{\mc{R}}^0_{\Phi}[F](y')
  -
  \partial^\gamma\dot{\mc{R}}^0_{\Phi}[F](y)
  \right|
  \\
   & \qquad\le
  A_\infty
  \int_0^\infty
  t^\alpha
  \int_{\R^d}
  \left|
  \partial_y^\gamma g_{x,t}(y')
  -
  \partial_y^\gamma g_{x,t}(y)
  \right|
  \dd x
  \frac{\dd t}{t}
  \\
   & \qquad\lesssim
  A_\infty
  \int_0^\infty
  t^{\alpha-n}
  \min\left(1,\frac{h}{t}\right)
  \frac{\dd t}{t}
  \lesssim
  A_\infty h^{\alpha-n}.
\end{aligned}
\]
Since \(\alpha-n=\delta\), this gives
\[
\left\|\dot{\mc{R}}^0_{\Phi}[F]\right\|_{\dot{\mathcal{C}}^\alpha(\R^d)}
\lesssim
A_\infty .
\]

We now prove the nonhomogeneous estimate.
Set
\[
A_0\eqd t_0^\alpha \left\| \1_{0<t<t_0}t^{-\alpha}F_0 \right\|_{L^\infty(\R^{d+1}_+)},
\qquad A_1\eqd \|F_1\|_{L^\infty(\R^d)}.
\]
Then \(A_{t_0}=A_0+A_1\) and
\[
|F_0(x,t)|\le A_0t_0^{-\alpha}t^\alpha, \qquad 0<t<t_0, \qquad |F_1(x)|\le A_1.
\]
Let \(|\beta|\le n\).
By \eqref{eq:common-wave-packet-l1-bound},
\[
\begin{aligned}
\left\|\partial^\beta \int_0^{t_0}\int_{\R^d} F_0(x,t) (\Tr_x\Dil_t\varphi_{x,t}) \dd x\frac{\dd t}{t} \right\|_{L^\infty}
& \lesssim A_0t_0^{-\alpha} \int_0^{t_0} t^{\alpha-|\beta|} \frac{\dd t}{t}
\\
& \lesssim A_0t_0^{-|\beta|},
\end{aligned}
\]
because \(|\beta|\le n<\alpha\).
Similarly,
\[
\left\| \partial^\beta \int_{\R^d} F_1(x) (\Tr_x\Dil_{t_0}\varphi^0_x) \dd x \right\|_{L^\infty} \lesssim A_1t_0^{-|\beta|}.
\]
Thus
\[
t_0^{|\beta|}\|\partial^\beta G\|_{L^\infty(\R^d)} \lesssim A_{t_0}, \qquad |\beta|\le n.
\]
The same estimates also give absolute pointwise convergence and locally uniform convergence after applying \(\partial^\beta\), for every \(|\beta|\le n\).

Finally, let \(|\gamma|=n\) and \(0<h=|y'-y|<t_0\).
By \eqref{eq:common-wave-packet-difference-bound},
\[
\begin{aligned}
& \left| \partial^\gamma \int_0^{t_0}\int_{\R^d} F_0(x,t) (\Tr_x\Dil_t\varphi_{x,t})(y') \dd x\frac{\dd t}{t}
- \partial^\gamma \int_0^{t_0}\int_{\R^d} F_0(x,t) (\Tr_x\Dil_t\varphi_{x,t})(y) \dd x\frac{\dd t}{t} \right|
\\
& \qquad\lesssim A_0t_0^{-\alpha} \int_0^{t_0} t^{\alpha-n} \min\left(1,\frac{h}{t}\right) \frac{\dd t}{t}
\lesssim A_0t_0^{-\alpha}h^{\alpha-n}.
\end{aligned}
\]
For the top-scale term, again using \eqref{eq:common-wave-packet-difference-bound},
\[
\begin{aligned}
& \left| \partial^\gamma \int_{\R^d} F_1(x) (\Tr_x\Dil_{t_0}\varphi^0_x)(y') \dd x - \partial^\gamma \int_{\R^d} F_1(x) (\Tr_x\Dil_{t_0}\varphi^0_x)(y) \dd x \right|
\\
& \qquad\lesssim A_1t_0^{-n} \min\left(1,\frac{h}{t_0}\right) \lesssim A_1t_0^{-\alpha}h^{\alpha-n},
\end{aligned}
\]
because \(0<h<t_0\) and \(0<\alpha-n<1\).
Therefore
\[
t_0^\alpha \sup_{\substack{y\ne y'\\ |y'-y|<t_0}} \frac{ |\partial^\gamma G(y')-\partial^\gamma G(y)| }{|y'-y|^{\alpha-n} }
\lesssim A_{t_0}, \qquad |\gamma|=n.
\]
Together with the scaled \(L^\infty\)-bounds for derivatives of order at most \(n\), this proves
\[
\|G\|_{\mathcal{C}^{\alpha}_{t_0}(\R^d)} \lesssim A_{t_0}.
\]
\end{proof}
